\documentclass[11pt]{article}
\usepackage[margin=1in]{geometry}
\usepackage[T1]{fontenc}
\usepackage{amsmath,amssymb}
\usepackage{booktabs}
\usepackage{enumitem}
\usepackage{hyperref}
\usepackage{xcolor}

\newcommand{\code}[1]{\texttt{\detokenize{#1}}}
\pdfstringdefDisableCommands{\def\code#1{#1}}
\newcommand{\h}[1]{\hat{#1}}
\newcommand{\hne}{\hat n_e}
\newcommand{\nFC}{\hat n_{\mathrm{FC}}}
\newcommand{\nCX}{\hat n_{\mathrm{CX}}}
\newcommand{\SCX}{S_{\mathrm{CX}}}
\newcommand{\fFC}{f_{\mathrm{FC}}}
\newcommand{\fCX}{f_{\mathrm{CX}}}
\newcommand{\VFC}{V_{\mathrm{FC}}}
\newcommand{\VCX}{V_{\mathrm{CX}}}
\newcommand{\gradr}{\langle|\nabla r|^2\rangle}

\title{\code{src/solver_3d.py}: Documentation of the Non-Dimensional Coupled\\
Three-Equation Saarelma--Connor Pedestal Solver (``3D'' model)}
\author{}
\date{}

\begin{document}
\maketitle
\tableofcontents

%=====================================================================
\section{Overview}
%=====================================================================
\code{src/solver_3d.py} is the non-dimensional solver for the coupled
three-equation Saarelma--Connor pedestal model (electron density plus the
Franck--Condon (FC) and charge-exchange (CX) neutral densities). It is
available in a Firedrake (finite-element) and a scipy (\code{solve_bvp}
collocation) discretisation. The module contributes
\code{ThreeDSolverMixin} to the single public class
\code{src.saarelma_connor_api.saarelma_connor}; all equilibrium, kinetic
and cross-section setup comes from \code{SaarelmaConnorBase}.

\paragraph{Entry point.} \code{saarelma_connor.solve_coupled_nondim(...)},
or equivalently \code{saarelma_connor.solve(model='3D', ...)}. Both return
the common result dictionary described in
\code{SaarelmaConnorBase._build_result_dict}.

\paragraph{Free parameters.} $\alpha_{\mathrm{crit}}$, \code{De_chie_etg},
$C_{\mathrm{KBM}}$, \code{nFC_x0}; set through the parent constructor,
through \code{update_free_params}, or per solve through the
\code{free_params} argument.

%=====================================================================
\section{Non-dimensionalisation}
%=====================================================================
The solver implements the foundational non-dimensionalisation derived in
Appendix~A.8 of \code{docs/Labbate_APAM9301_Final_06152026.tex}
(Eq.~(eq:nondim-scales-A8)):
\begin{align*}
L &= |x_{\mathrm{inner}}| && \text{length scale (m)},\\
n_0 &= \code{ne_x0} && \text{density scale (m}^{-3}\text{), separatrix value},\\
T_0 &= T_e(0)+T_i(0) && \text{energy scale (J), separatrix value},\\
\tau &= 1/(n_0 S_0),\quad S_0 = S_i(0) && \text{time scale (s)},
\end{align*}
with derived scales
\[
[V]_0 = L/\tau = L n_0 S_0 \;(\mathrm{m/s}),\qquad
[D]_0 = L^2/\tau = L^2 n_0 S_0 \;(\mathrm{m^2/s}),\qquad
[p]_0 = n_0 T_0 \;(\mathrm{Pa}).
\]

\paragraph{Hatted variables.}
\begin{align*}
\h x &= x/L,\ \h x\in[-1,0], &
\hne &= n_e/n_0,\ \hne(0)=1, \\
\nFC &= \langle n_{\mathrm{FC}}\rangle/n_0,\ \nFC(0)=\code{nFC_x0}/n_0, &
\nCX &= \langle n_{\mathrm{CX}}\rangle/n_0,\ \nCX(0)=\code{nCX_x0}/n_0,\\
\h T &= (T_e+T_i)/T_0, &
\h S_i &= S_i/S_0,\quad \h S_{\mathrm{CX}} = \SCX/S_0,\\
\h V_{\mathrm{FC}} &= |\VFC|/[V]_0,\quad \h V_{\mathrm{CX}} = |\VCX|/[V]_0, &
\h g &= \gradr \text{ (already dimensionless)}.
\end{align*}

\paragraph{Dimensionless transport coefficients.}
\[
\h C_{\mathrm{ETG}} = \frac{C_{\mathrm{ETG}}}{L^2 n_0^2 S_0},\quad
\h D_{\mathrm{NEO}} = \frac{D_{\mathrm{NEO}}}{L^2 n_0 S_0},\quad
\h A_{\mathrm{KBM}} = \frac{A_{\mathrm{KBM}}}{L^2 n_0 S_0},\quad
\h B_{\mathrm{KBM}} = \frac{B_{\mathrm{KBM}}\,T_0}{L^3 S_0}.
\]
For readability, two composite scales are cached:
\code{_ETG_scale_nd} $= L^2 n_0^2 S_0 = n_0[D]_0$ and
\code{_KBMB_scale_nd} $= L^3 S_0/T_0 = [B_{\mathrm{KBM}}]$.

%=====================================================================
\section{Dimensionless residuals}
%=====================================================================
The residuals are Eqs.~(eq:weak-hat-A8), (eq:weak-hat-nFC-A8),
(eq:weak-hat-nCX-A8) of App.~A.8, with test functions $v_e$, $v_F$, $v_C$:
\begin{align}
\h F_1 &= \int_{-1}^{0} \h g\Big[\frac{\h C_{\mathrm{ETG}}}{\hne}\hne'
   + \h A_{\mathrm{KBM}}\hne' + \h B_{\mathrm{KBM}}\h T(\hne')^2
   + \h B_{\mathrm{KBM}}\hne\,\frac{d\h T}{d\h x}\,\hne'
   + \h D_{\mathrm{NEO}}\hne'\Big] v_e'\,d\h x \nonumber\\
 &\quad - \int_{-1}^{0}\hne\h S_i(\nFC+\nCX)\,v_e\,d\h x
 + \delta_N\,\h g(-1)\,\h D|_{-1}\,\h{n}'_{e,\mathrm{inner}}\,v_e(-1) = 0,
 \label{eq:F1}\\
\h F_2 &= \int_{-1}^{0}\Big[\h V_{\mathrm{FC}}\frac{d}{d\h x}\bigl[\fFC\h g\nFC\bigr]
   - \hne(\h S_i+\h S_{\mathrm{CX}})\nFC\Big] v_F\,d\h x = 0,
 \label{eq:F2}\\
\h F_3 &= \int_{-1}^{0}\Big[\frac{d}{d\h x}\bigl[\h V_{\mathrm{CX}}\fCX\h g\nCX\bigr]
   - \hne\bigl(\h S_i\nCX - \tfrac12\h S_{\mathrm{CX}}\nFC\bigr)\Big] v_C\,d\h x = 0.
 \label{eq:F3}
\end{align}

%=====================================================================
\section{Placement of the $n_e$ gradient boundary condition}
\label{sec:bcs}
%=====================================================================
\code{ne_grad_bc_loc} selects which end of the domain carries the
prescribed $dn_e/dx$.

\paragraph{\code{"inner"} (default, unchanged behaviour).} One condition at
each end:
\begin{align*}
\hne(0) &= 1, \qquad \nFC(0) = \code{nFC_x0}/n_0, \qquad \nCX(0) = \code{nCX_x0}/n_0,\\
\hne'(-1) &= L\,\code{dne_dx_inner}/n_0 \quad\text{(Neumann flux on ds(1))}.
\end{align*}
(A Dirichlet inner condition $\hne(-1)=\code{ne_inner}/n_0$ was the legacy
alternative; \code{ne_inner_bc} values other than \code{"neumann"} are now
rejected by \code{bc_ig_helpers.reject_legacy_ne_inner_bc}.)

\paragraph{\code{"outer"}.} Both $n_e$ conditions sit at the separatrix,
and the inner boundary is left completely free:
\[
\hne(0) = 1,\quad \hne'(0) = L\,\code{dne_dx_outer}/n_0,\quad
\nFC(0) = \code{nFC_x0}/n_0,\quad \nCX(0) = \code{nCX_x0}/n_0,
\]
with nothing imposed at $\h x=-1$. This is the boundary-condition set of
Saarelma \emph{et al.} 2023 Sec.~2.3, where the separatrix density and its
gradient (their Eq.~(20), $dn_e/dx|_0 = -n_e(0)/\sqrt{D_{\mathrm{SOL}}\tau_\parallel}$)
are the specified data and nothing is imposed at the pedestal top.

A Galerkin discretisation cannot impose both directly: the strong
Dirichlet condition at boundary id 2 makes $v_e(0)=0$, so a \code{ds(2)}
flux term vanishes identically. The equivalent well-posed statement is
that the \emph{inner} flux is an unknown Lagrange multiplier fixed by the
separatrix-gradient constraint, and that is what is solved here --- by
secant iteration on the \code{ds(1)} slope \code{hat_dne_dx_inner} until
$\hne'(0)$ matches the prescribed value (\code{_shoot_outer_grad_nondim},
Section~\ref{sec:shoot}). Shooting on that one scalar (rather than
enlarging the mixed space by a real-space multiplier) leaves the function
space, the Picard loop and the analytic-neutral path untouched. Diagnostics
land in \code{self.grad_bc_info}.

The scipy discretisation needs no shooting: \code{solve_bvp} only needs
the right \emph{number} of residuals and does not care which end they are
evaluated at.

Boundary-condition resolution and initial guesses live in
\code{src/bc_ig_helpers.py} so the 1D/3D $\times$ Firedrake/scipy solvers
share exactly one implementation.

%=====================================================================
\section{Class \code{ThreeDSolverMixin}}
%=====================================================================
Non-dimensional coupled three-equation solver (``3D'' model). A mixin over
\code{SaarelmaConnorBase}; it is combined with the base and
\code{OneDSolverMixin} into the single public \code{saarelma_connor}. It
defines no \code{__init__} --- every attribute it reads is set by the base
constructor --- and it overrides nothing, so it only ever adds the
\code{*_nondim} methods. The solver assembles and solves the dimensionless
residuals of App.~A.8 rather than their SI counterparts.

Reference scales are computed lazily inside the solve (after
\code{setup_solver_grids} and any inner-boundary search are done) and
stored as instance attributes:
\begin{itemize}[nosep]
  \item \code{_L_nd}, \code{_n0_nd}, \code{_T0_nd}, \code{_S0_nd}: the four free scales;
  \item \code{_tau_nd}, \code{_V0_nd}, \code{_D0_nd}: derived scales.
\end{itemize}
After convergence the SI-unit profiles are recovered and stored on the same
attributes as the parent solver (\code{x_sol}, \code{ne_sol},
\code{nFC_sol}, \code{nCX_sol}), so all downstream code (plotting, EPED
feed, etc.) keeps working without modification. The non-dimensional
profiles are also kept as \code{hat_x_sol}, \code{hat_ne_sol},
\code{hat_nFC_sol}, \code{hat_nCX_sol} for diagnostics.

%=====================================================================
\section{Method reference}
%=====================================================================

\subsection{\code{_set_nondim_scales(verbose=False)}}
Compute and cache the four free reference scales of
Eq.~(eq:nondim-scales-A8) plus the derived scales:
\begin{center}
\begin{tabular}{lll}
\toprule
Scale & Definition & Units \\
\midrule
$L$ & $|x_{\mathrm{inner}}|$ & m \\
$n_0$ & \code{ne_x0} (separatrix value) & m$^{-3}$ \\
$T_0$ & $(T_e(0)+T_i(0))\cdot$ J/eV (separatrix value) & J \\
$S_0$ & $S_i(0)$ (separatrix value) & m$^3$/s \\
$\tau$ & $1/(n_0 S_0)$ & s \\
$[V]_0$ & $L/\tau$ & m/s \\
$[D]_0$ & $L^2/\tau$ & m$^2$/s \\
\bottomrule
\end{tabular}
\end{center}
\code{T_e_pres} / \code{T_i_pres} are stored in eV, and the separatrix is
index $-1$ because \code{x_init = r_psi - r_psi[-1]}; $S_i$ at the
separatrix is taken from \code{S_i_pres} (built in
\code{setup_solver_grids} on \code{psi_N_pres}). Raises if
\code{x_inner} / \code{ne_x0} are unset or if $S_i(0)$ is not finite and
positive.

\subsection{\code{_to_hat_x(x_si)}, \code{_from_hat_x(hat_x)}}
Map an SI $x$ array (m) onto $\h x = x/L$, and back ($x = L\h x$).

\subsection{\code{_interp_si_to_hat(hat_x_dofs, si_array_on_xinit)}}
Evaluate an SI quantity defined on \code{self.x_init} (sorted) at the
dimensionless mesh DOFs \code{hat_x_dofs}. Returns the SI value at
$L\cdot\code{hat_x_dofs}$ (i.e.\ \emph{no} rescaling); callers rescale to
hat-units as needed. In effect this just interpolates the SI array onto
the \code{x_dofs} grid; because \code{x_dofs} has the same size as
\code{hat_x_dofs}, no further interpolation is needed once there.

\subsection{\code{_ensure_firedrake_discretization_nondim(mesh_n, fe_degree, force=False)}}
Build (or reuse) the dimensionless mesh on $\h x\in[-1,0]$ (CG of degree
\code{fe_degree}, plus the three-field mixed space $W = V\times V\times V$)
and the equilibrium-only non-dimensional coefficient Functions
$\h g$, $\h S_i$, $\h S_{\mathrm{CX}}$, $\h V_{\mathrm{CX}}$. Cache keys
live alongside the parent's \code{self._fd_cache} keys but use a
\code{_nd} suffix to avoid colliding with any cached SI-version state; any
mixed solution cached against a previous mesh is dropped.

Returns \code{(mesh, V, W, hat_x_dofs, hat_g_fd, hat_Si_fd, hat_Scx_fd, hat_Vcx_fd)}.

\subsection{\code{_get_or_create_mixed_solution_nondim(W, force=False)}}
Return cached \code{(hat_u, hat_u_prev)} mixed Functions or allocate them.

\subsection{\code{calc_pressure_quantities_nondim(hat_n_e, gate_mode=None, hat_x_dofs=None)}}
Non-dimensional version of \code{calc_pressure_quantities}. Computes the
Connor--Hastie $\alpha$ in physical (SI) units (because $\alpha$ is a
physical quantity defined by an SI integral),
\[
\alpha = \alpha_{\mathrm{nodp}}\,\frac{dp}{dx},\qquad p = n_e(T_e+T_i)e
\]
(neutral pressures neglected, quasi-neutral plasma assumed), and then
rescales the resulting transport contributions $A_{\mathrm{KBM}}$,
$B_{\mathrm{KBM}}$, $D_{\mathrm{KBM}}$ into their dimensionless
counterparts (App.~A.8 Eqs.~(eq:hat-A-KBM), (eq:hat-B-KBM)):
\[
\h A = A/[D]_0,\qquad \h B = B\,T_0/(L^3 S_0) = B/\code{KBMB_scale},\qquad
\h D_{\mathrm{KBM}} = D_{\mathrm{KBM}}/[D]_0.
\]
With $G_{\mathrm{KBM}} = C_{\mathrm{KBM}}c_s\rho_s^2/a$, the SI
coefficients are $A = -G_{\mathrm{KBM}}\alpha_{\mathrm{crit}}$,
$B = G_{\mathrm{KBM}}\alpha_{\mathrm{nodp}}$ where the gate is on.

\paragraph{Parameters.}
\begin{description}[style=nextline]
  \item[\code{hat_n_e} (array, shape (n_dofs,))] Dimensionless electron
  density at the mesh DOFs (DOF order, not necessarily monotonic in $\h x$).
  \item[\code{gate_mode} (\code{"average"} | \code{"majority"} | \code{"local"})]
  How the KBM on/off gate is decided:
  \begin{itemize}[nosep]
    \item \code{"average"}: gate on the pedestal-averaged $\alpha$,
    $\text{gate} = \mathrm{mean}(\alpha) > \alpha_{\mathrm{crit}}$ (single
    on/off for the whole grid). When on, Saarelma \emph{et al.} (2023)
    Eqs.~24--25 diffusivity is used on the whole grid:
    $D_{\mathrm{KBM}} = (\bar\alpha - \alpha_{\mathrm{crit}})G_{\mathrm{KBM}}$.
    $A/B$ are kept only as diagnostics (the Picard ``average'' path freezes
    $D$ into the $A$-slot with $B=0$).
    \item \code{"majority"}: gate on a majority vote of the local $\alpha$:
    on for the whole grid iff more than half of the pedestal grid points have
    $\alpha > \alpha_{\mathrm{crit}}$. When on, the local $(A,B)$ KBM
    coefficient structure is applied at every point (no pointwise gate, and
    no global $\bar\alpha$ lock as in \code{"average"}).
    \item \code{"local"}: pointwise gate (legacy behaviour): each grid point
    is gated by its own local $\alpha$.
  \end{itemize}
  \item[\code{hat_x_dofs} (array or None)] Grid on which \code{hat_n_e} is
  given. The Firedrake path leaves this \code{None} and uses the cached FE
  DOFs; the scipy path passes its own (ascending) collocation grid.
\end{description}

\paragraph{Sets.} \code{_hat_A_KBM}, \code{_hat_B_KBM}, \code{_hat_D_KBM}
(dimensionless KBM coefficients in DOF order); \code{alpha_bar_ped}
(pedestal-averaged $\alpha$); \code{alpha_frac_above} (fraction of pedestal
points with $\alpha>\alpha_{\mathrm{crit}}$); \code{kbm_gate_on} (scalar
bool for \code{"average"}/\code{"majority"}, pointwise bool array in DOF
order for \code{"local"}); \code{alpha_local_ped} (local $\alpha$ in DOF
order); and the diagnostics \code{D_KBM_si}, \code{G_KBM},
\code{alpha_nodp}, \code{pres}, \code{T_e_xdofs}.

\subsection{\code{_solve_neutrals_analytic_nondim(hat_ne_dofs, n_sub=4001)}}
Solve the FC and CX neutral equations exactly for a \emph{given} electron
density, on a dense sub-grid that resolves the neutral ionisation boundary
layer even when the FE mesh does not.

Both neutral equations are first-order linear ODEs in $x$ once $n_e(x)$ is
frozen, so they admit integrating-factor solutions that are positive by
construction (no CG oscillations, no negative densities). This is the same
mathematics as the \code{nFC_ic='solve'} / \code{nCX_ic='solve'}
initial-guess branches, but evaluated on \code{n_sub} points instead of the
FE DOFs and including the $\gradr$ / form-factor variation exactly.

\paragraph{FC neutrals} (Saarelma--Connor Eq.~9):
\[
|\VFC|\frac{d}{dx}\bigl[\fFC\,g\,n_{\mathrm{FC}}\bigr] = n_e(S_i+\SCX)\,n_{\mathrm{FC}}.
\]
With $u = \fFC g n_{\mathrm{FC}}$ and $\tau = -x$ (inward distance $\ge0$):
\[
u(\tau) = u(0)\exp\!\left(-\int_0^\tau \frac{n_e(S_i+\SCX)}{|\VFC|\fFC g}\,d\tau'\right).
\]

\paragraph{CX neutrals} (Eq.~10), with $w = |\VCX|\fCX g n_{\mathrm{CX}}$:
\[
\frac{dw}{d\tau} = -P(\tau)w + Q(\tau),\qquad
P = \frac{n_e S_i}{|\VCX|\fCX g},\qquad
Q = \tfrac12 n_e \SCX n_{\mathrm{FC}},
\]
integrated with a per-step exponential integrator
\[
w_{k+1} = w_k e^{-P\Delta\tau} + \frac{Q}{P}\bigl(1-e^{-P\Delta\tau}\bigr)
\]
(using midpoint $P$, $Q$; falls back to $w_{k+1}=w_k+Q\Delta\tau$ when
$P\Delta\tau < 10^{-12}$), which is unconditionally stable and
positivity-preserving regardless of how stiff $P\Delta\tau$ is.

\paragraph{Parameters.}
\begin{description}[style=nextline]
  \item[\code{hat_ne_dofs} (array)] Dimensionless $n_e$ at the FE mesh
  DOFs (DOF order).
  \item[\code{n_sub} (int, 4001)] Number of points of the dense integration
  sub-grid on $[x_{\mathrm{inner}},0]$. Choose so that the sub-grid spacing
  is well below the minimum neutral penetration depth
  $\lambda = |\VFC|/(n_e(S_i+\SCX))$.
\end{description}

\paragraph{Returns} \code{(hat_nFC_dofs, hat_nCX_dofs)}: dimensionless
neutral densities at the FE mesh DOFs (DOF order), interpolated in
log-space from the sub-grid so the boundary-layer decay is preserved. The
SI sub-grid profiles are also stored as \code{x_neutrals_analytic},
\code{nFC_analytic}, \code{nCX_analytic}.

\subsection{\code{_kbm_inline_terms(...)}}
Build the inline (trial-dependent) UFL expressions for the Connor--Hastie
$\alpha$, the smoothed KBM gate, and the dimensionless KBM coefficients
$\h A_{\mathrm{KBM}}$ and $\h B_{\mathrm{KBM}}$ (App.~A.7--A.8):
\begin{align*}
\alpha(\h x,\hne) &= \h\alpha_{\mathrm{nodp}}(\h x)\,\frac{d}{d\h x}\bigl[\hne\h T\bigr]
 = \h\alpha_{\mathrm{nodp}}\bigl(\h T\hne' + \hne\h T'\bigr),\\
\mathrm{gate}_\epsilon(z) &= \tfrac12\bigl(1+\tanh(z/\epsilon)\bigr) \approx H(z),\\
\h A_{\mathrm{KBM}}(\hne) &= -\mathrm{gate}_\epsilon(\alpha-\alpha_{\mathrm{crit}})\,\alpha_{\mathrm{crit}}\,\h G_{\mathrm{KBM}},\\
\h B_{\mathrm{KBM}}(\hne) &= +\mathrm{gate}_\epsilon(\alpha-\alpha_{\mathrm{crit}})\,\h\alpha_{\mathrm{nodp}}\,\h G_{\mathrm{KBM}}.
\end{align*}
These are full UFL expressions, so SNES (Newton) automatically captures
their dependence on the trial $\hne$.

\code{boundary} (bool, default \code{False}): if \code{True}, use the
prescribed inner-boundary slope \code{hat_dne_dx_inner_c} for $\hne'$
inside $\alpha$ (needed in the Neumann-flux \code{ds(1)} contribution to
$F_1$, so the gate at the inner boundary stays consistent with the imposed
flux). Otherwise the volume trial derivative \code{hat_ne.dx(0)} is used.

Returns \code{(alpha_ufl, gate_ufl, hat_A_KBM_ufl, hat_B_KBM_ufl)}.

\subsection{\code{_build_f1_weak_form_nondim(...)}}
Dimensionless $n_e$ weak form, Eq.~\eqref{eq:F1} (Eq.~(eq:weak-hat-A8) of
App.~A.8). Integrals are over $\h x\in[-1,0]$ (the Firedrake mesh). The
Neumann boundary contribution at $\h x=-1$ (boundary id 1) is always
present --- the prescribed inner flux in the \code{"inner"} pathway, the
free (secant-driven) inner flux in the \code{"outer"} pathway --- evaluated
using the same expansion of $\h D$ as the volume term, with the prescribed
slope substituted for $\hne'$:
\[
\h D_{bc} = \frac{\h C_{\mathrm{ETG}}}{\hne} + \h A^{bc}_{\mathrm{KBM}}
 + \h B^{bc}_{\mathrm{KBM}}\h T\,\hne'_{bc} + \h B^{bc}_{\mathrm{KBM}}\hne\h T'
 + \h D_{\mathrm{NEO}},\qquad
F_1 \mathrel{+}= \h g\,\h D_{bc}\,\hne'_{bc}\,v_e\,ds(1).
\]
\code{hat_A_KBM_term} and \code{hat_B_KBM_term} are UFL expressions (or
Functions) carrying the KBM coefficients; in the \code{"inline"} path they
depend on the trial $\hne$. \code{hat_A_KBM_bc_term} /
\code{hat_B_KBM_bc_term} are the analogous expressions evaluated with the
prescribed Neumann slope (used only inside the \code{ds(1)} integrand); if
left \code{None} the volume expressions are reused.

\subsection{\code{_build_f2_weak_form_nondim(...)}}
Dimensionless FC-neutral weak form, Eq.~\eqref{eq:F2}
(Eq.~(eq:weak-hat-nFC-A8)).

\subsection{\code{_build_f3_weak_form_nondim(...)}}
Dimensionless CX-neutral weak form, Eq.~\eqref{eq:F3}
(Eq.~(eq:weak-hat-nCX-A8)).

\subsection{\code{_shoot_outer_grad_nondim(F, u, bcs, snes_params, hat_ne_curr, hat_slope_c, hat_target, tol, max_it, verbose)}}
\label{sec:shoot}
Solve $F=0$ subject to $\hne'(0) = \code{hat_target}$.

The inner boundary carries no condition of its own in this mode, so its
flux slope \code{hat_slope_c} (the \code{ds(1)} Constant, which the
inline-KBM boundary terms share) is the free unknown. A damped secant
iteration drives the residual
\[
r(s) = \hne'(0;s) - \code{hat_target}
\]
to zero ($\hne'(0)$ is assembled as \code{hat_ne.dx(0)*ds(2)}, since
\code{ds(2)} has unit measure in 1D; \code{tol} is relative to
$|\code{hat_target}|$).

\paragraph{Step control.} The seed is the inner slope, which can be far from
the answer (the free flux routinely lands on the other side of zero), and
the system need not be solvable for every slope in between, so each step is
capped at four times the previous accepted one and backtracked, halving (up
to 12 times), whenever SNES fails. The first step is a finite-difference
bootstrap of $0.05|s|$ (or $10^{-3}$ if $s=0$).

\paragraph{Fixed restart state.} Every attempt restarts from the
\emph{same} reference state (the iterate this routine was handed) rather
than from the previous attempt's answer. That is deliberate. Chaining warm
starts makes $r(s)$ path dependent, and once two slopes are close the warm
start already passes SNES's convergence test, so it returns at iteration
zero with the density untouched; the secant then sees a bit-identical
residual for two different slopes and concludes the gradient does not
respond at all. Restarting from a fixed reference makes $r(s)$ a genuine
function of $s$, which is what the secant needs; the cost is a few more
Newton steps per attempt. \code{snes_stol} is set to 0 so SNES cannot exit
on the step-size test before doing any work.

\paragraph{Returns} a diagnostics dict (\code{converged}, \code{iterations},
\code{solves}, \code{backtracks}, \code{hat_slope}, \code{rel_residual},
\code{history}), also stored by the caller in \code{self.grad_bc_info}.
Raises \code{RuntimeError} on a failed seed solve, a stalled secant,
repeated SNES failure, or exhaustion of \code{max_it}.

\subsection{\code{_check_solver_structure(solver_structure)} (static)}
Validate and normalise the discretisation-choice flag (\code{"firedrake"}
or \code{"scipy"}).

\subsection{\code{_scipy_coefficients_nondim()}}
Linear interpolants of the frozen non-dimensional coefficients in $\h x$
(keys \code{g}, \code{Si}, \code{Scx}, \code{Vcx}, \code{fFC}, \code{fCX},
\code{D_NEO}, \code{C_ETG}, and the scalar \code{VFC}). Every coefficient
lives on \code{self.x_init} in SI; each returned callable takes $\h x$ and
applies $x = L\h x$ internally, so the ODE right-hand side never has to
think about units. Linear interpolation matches the \code{np.interp}
sampling the Firedrake path uses to fill its coefficient Functions, so the
two discretisations see the same coefficient reconstruction.

%---------------------------------------------------------------------
\subsection{\code{solve_coupled_nondim_scipy(...)}}
\label{sec:scipy}
%---------------------------------------------------------------------
scipy (\code{solve_bvp}) solver for the coupled three-equation model. Same
physics, non-dimensionalisation and boundary conditions as
\code{solve_coupled_nondim}, discretised by collocation instead of finite
elements. Reached through
\code{solve_coupled_nondim(solver_structure="scipy", ...)}.

\paragraph{Formulation.} Equations (8)--(10) of the writeup form a
\emph{fourth}-order system (second order in $\hne$, first order in each
neutral), which is exactly Saarelma \emph{et al.} 2023 Sec.~2.3's ``fourth
order system requiring four boundary conditions''. \code{solve_bvp} wants
an explicit first-order system, so the state carries the three
\textbf{fluxes} rather than the derivatives:
\[
Y = [\hne,\ \Phi,\ U,\ W],\qquad
\Phi = \h f\,\hne' \ \text{(electron flux)},\quad
U = \fFC\h g\nFC \ \text{(FC flux}/\h V_{\mathrm{FC}}),\quad
W = \h V_{\mathrm{CX}}\fCX\h g\nCX \ \text{(CX flux)},
\]
with the conductance
$\h f(\h x,\hne) = \h g(\h D_{\mathrm{NEO}}+\h D_{\mathrm{KBM}}) + \h g\h C_{\mathrm{ETG}}/\hne$.
The system is then
\begin{align*}
\hne' &= \Phi/\h f,\\
\Phi' &= -\hne\h S_i(\nFC+\nCX),\\
U' &= \hne(\h S_i+\h S_{\mathrm{CX}})\nFC/\h V_{\mathrm{FC}},\\
W' &= \hne\bigl(\h S_i\nCX - \h S_{\mathrm{CX}}\nFC/2\bigr),
\end{align*}
recovering $\nFC = U/(\fFC\h g)$ and $\nCX = W/(\h V_{\mathrm{CX}}\fCX\h g)$.
Writing it this way means \textbf{no coefficient derivatives appear
anywhere} --- unlike the 1D \code{solve_sc_scipy}, which expands
$d/dx[f n_e']$ and therefore needs $df/dx$ off the coarse p-file grid (the
accuracy-limiting step of that solver). Here the divergence form is kept
intact, as in the Firedrake weak form.

\paragraph{Boundary conditions.} Four residuals, placed exactly as in the
Firedrake path (Section~\ref{sec:bcs}): $\hne(0)$, $U(0)$ and $W(0)$ are
pinned at the separatrix; the fourth is $\Phi(0) = \h f(0,\hne(0))\,\hne'_{bc}$
in \code{"outer"} mode or $\Phi(-1) = \h f(-1,\hne(-1))\,\hne'_{bc}$ in
\code{"inner"} mode. \code{solve_bvp} only needs the right \emph{number}
of residuals, so \code{"outer"} needs no shooting here.

\paragraph{Initial guess.} $n_e$ comes from
\code{bc_ig_helpers.build_ne_initial_guess} (\code{initial_guess},
\code{tanh_width}, \code{tanh_center}) and the two neutral profiles from
\code{bc_ig_helpers.build_neutral_initial_guess} (\code{nFC_ic},
\code{nCX_ic}) --- the same shared helpers, with the same options, that the
Firedrake path uses, so the two discretisations start from identical SI
profiles. \code{nFC_ic="solve"} / \code{nCX_ic="solve"} integrate the
neutral equations analytically at the frozen $n_e$ guess;
\code{nFC_ic} / \code{nCX_ic} \code{="manual EPEDNN loop"} read the
tabulated EPED-NN profiles, and \code{nCX_ic="scale nFC"} scales the FC
guess by $\code{nCX_x0}/\code{nFC_x0}$. The SI neutral densities are
converted to the flux variables $U$ and $W$ before being handed to
\code{solve_bvp}.

\paragraph{KBM treatment.} Only \code{picard_gate_mode="average"} is
supported: \code{"majority"} freezes the local $A/B$ KBM structure, which
makes the conductance depend on $\hne'$ and turns $\Phi\to\hne'$ into a
per-point root-find; use \code{solver_structure="firedrake"} for it.
$D_{\mathrm{KBM}}$ is frozen per Picard iteration (initially from the
initial guess, Saarelma Eq.~25) and enters the ODE through an interpolant
rebuilt each time round the loop, with optional under-relaxation
\code{picard_relax}. \code{ne_floor} clips $\hne$ from below so that
$\h C_{\mathrm{ETG}}/\hne$ stays finite.

\paragraph{Returns.} The same result schema as \code{solve_coupled_nondim},
in SI units. $dn_e/dx$ is recovered as $(\Phi/\h f)(n_0/L)$.

%---------------------------------------------------------------------
\subsection{\code{solve_coupled_nondim(...)}}
\label{sec:driver}
%---------------------------------------------------------------------
Non-dimensional Firedrake solver for the coupled three-equation
Saarelma--Connor neutral-transport pedestal model. Equivalent to
\code{saarelma_connor.solve_coupled} but assembles and solves the
dimensionless residuals derived in App.~A.8 (foundational scaling, no free
$\Lambda$ prefactor). Inputs and outputs are in SI units; the rescaling is
invisible to the caller. With \code{solver_structure="scipy"} it forwards
to \code{solve_coupled_nondim_scipy}.

\paragraph{Boundary conditions.} See Section~\ref{sec:bcs}. An inner
boundary location \code{psi_N_inner_boundary} must be specified (the
automatic inner-boundary search is no longer offered).

\paragraph{Parameters.} All physical inputs are in SI units.
\begin{description}[style=nextline]
  \item[\code{x_res} (int)] Number of mesh cells (Firedrake) or collocation
  points (scipy).
  \item[\code{free_params} (dict or None)] Optional
  \code{{'alpha_crit', 'De_chie_etg', 'C_KBM', 'nFC_x0'}} (plus any
  keyword accepted by \code{update_free_params}, e.g.\
  \code{psi_N_inner_boundary}) applied before the setup, exactly as in
  \code{solve_sc_firedrake}. Left \code{None} the four parameters must
  already be set on the instance; either way they are validated before
  anything is assembled. Forwarded unchanged when
  \code{solver_structure="scipy"}.
  \item[\code{fe_degree} (int, 2)] CG element degree.
  \item[\code{solver_structure} (\code{"firedrake"} | \code{"scipy"})]
  Which discretisation solves the coupled system. \code{"scipy"} forwards
  to \code{solve_coupled_nondim_scipy} (collocation via \code{solve_bvp});
  it supports \code{picard_gate_mode="average"} only. Firedrake-only
  arguments (\code{fe_degree}, \code{linear_solver}, \code{ksp_*},
  \code{kbm_treatment="inline"}, \code{neutrals_treatment},
  \code{grad_bc_*}) do not apply there.
  \item[\code{ne_inner_bc}] Must be \code{"neumann"} (legacy Dirichlet
  option rejected).
  \item[\code{ne_grad_bc_loc} (\code{"inner"} | \code{"outer"})] Which end
  of the domain carries the prescribed $dn_e/dx$. \code{"outer"} frees the
  inner boundary entirely (the \code{ds(1)} flux becomes the unknown that
  the separatrix gradient determines).
  \item[\code{bc_origin}, \code{ne_inner}, \code{dne_dx_inner}] Origin and
  values of the boundary data (resolved by
  \code{bc_ig_helpers.resolve_ne_bcs}).
  \item[\code{dne_dx_outer} (float or None)] SI separatrix slope
  $dn_e/dx|_{x=0}$ (m$^{-4}$) used when \code{ne_grad_bc_loc="outer"}.
  Taken from the same origins as the inner slope: given explicitly it wins
  for any \code{bc_origin} (this is how Saarelma Eq.~(20) is fed in); left
  \code{None} it is read off the p-file gradient at $x=0$, which requires a
  p-file-backed \code{bc_origin}.
  \item[\code{grad_bc_tol} (float, $10^{-8}$)] Relative tolerance on
  $\hne'(0)-\text{target}$ for the \code{"outer"} secant iteration.
  \item[\code{grad_bc_max_it} (int, 25)] Maximum secant iterations for the
  \code{"outer"} mode.
  \item[\code{grad_bc_seed} (float or None)] SI seed slope for the
  \code{"outer"} secant. Default is zero inner flux: the natural boundary
  condition a Galerkin form defaults to, so it is the neutral starting
  point and --- unlike seeding from a pedestal-top gradient --- uses no
  information the \code{"outer"} pathway forbids. It is also markedly more
  robust: $r(s)$ is discontinuous where the profile switches solution
  branch, and approaching from zero avoids straddling that jump.
  \item[\code{initial_guess}, \code{tanh_width}, \code{tanh_center}] Shape
  of the initial $n_e$ profile; \code{tanh_width} is in $x$ units (m).
  \item[\code{linear_solver}, \code{ksp_rtol}, \code{ksp_max_it}] PETSc
  linear-solver controls.
  \item[\code{reuse_setup} (bool)] Reuse cached grids / meshes.
  \item[\code{scale_ne_inner} (float or None)] Scaling factor for the inner
  boundary density, for testing purposes.
  \item[\code{nFC_ic}, \code{nCX_ic}] Neutral initial-guess options (see
  Section~\ref{sec:scipy}).
  \item[\code{bvp_tol}, \code{bvp_max_nodes}] \code{solve_bvp} controls,
  used only by \code{solver_structure="scipy"}.
  \item[\code{ne_floor} (float)] Lower clip on $\hne$ keeping
  $\h C_{\mathrm{ETG}}/\hne$ finite in the scipy driver.
  \item[\code{verbose} (bool or None)] Override \code{self.verbose}.
\end{description}

\paragraph{KBM treatment.} \code{kbm_treatment} controls how the piecewise
KBM gate (Heaviside on $\alpha-\alpha_{\mathrm{crit}}$) is handled.
\begin{itemize}
  \item \code{"inline"} (default): $\h A_{\mathrm{KBM}}$ and
  $\h B_{\mathrm{KBM}}$ are written directly as UFL expressions of the trial
  $\hne$ (App.~A.7--A.8), with the discontinuous indicator replaced by the
  smoothed Heaviside
  \[
  \mathrm{gate}(z) = \tfrac12\bigl(1+\tanh(z/\code{kbm_gate_eps})\bigr),
  \]
  so the residual is $C^\infty$ in $\hne$ and SNES Newton can differentiate
  it analytically. The Connor--Hastie $\alpha$ is evaluated locally as
  $\alpha(\h x,\hne) = \h\alpha_{\mathrm{nodp}}(\h x)\,d(\hne\h T)/d\h x$.
  The gate therefore updates self-consistently within every Newton step; on
  convergence the model is on its own KBM branch by construction. The
  required equilibrium fields are
  \[
  \h G_{\mathrm{KBM}} = \frac{C_{\mathrm{KBM}}c_s\rho_s^2/a}{L^2 n_0 S_0},\qquad
  \h\alpha_{\mathrm{nodp}} = \alpha_{\mathrm{nodp}}\,\frac{n_0 T_0}{L}.
  \]
  \item \code{"picard"}: outer Picard (fixed-point) loop following Saarelma
  \emph{et al.} (2023). Each Picard iteration the KBM gate and coefficients
  are evaluated from the \emph{current} $\hne$ (the initial guess on the first
  pass, thereafter the previously solved profile), then frozen while SNES
  solves the three-field system. The loop repeats until the density profile
  is unchanged to \code{picard_rtol} and the gate state is stable (or
  \code{picard_max_it} is hit). \code{picard_gate_mode} selects both the
  whole-grid gate criterion and the frozen flux structure:
  \begin{itemize}[nosep]
    \item \code{"average"} (default): KBM on iff the pedestal-averaged
    $\alpha$ exceeds $\alpha_{\mathrm{crit}}$ (Eq.~24). When on, freeze
    $\h D_{\mathrm{KBM}} = (\bar\alpha-\alpha_{\mathrm{crit}})\h G$ into the
    $A$-slot with $B=0$ (Eq.~25 diffusivity).
    \item \code{"majority"}: KBM on iff more than half of the pedestal grid
    points have local $\alpha>\alpha_{\mathrm{crit}}$. When on, freeze the
    local $A/B$ KBM structure everywhere (no pointwise gate).
  \end{itemize}
\end{itemize}
\begin{description}[style=nextline]
  \item[\code{kbm_gate_eps} (float or None)] Smoothing width of the
  Heaviside (inline treatment only). \code{None} $\to$ 5\% of
  $\alpha_{\mathrm{crit}}$ (with a $10^{-3}$ floor, avoiding division by
  zero at $\alpha_{\mathrm{crit}}=0$). Smaller $\epsilon$ makes the gate
  sharper but harder for Newton to converge.
  \item[\code{picard_max_it} (int, 50)] Maximum number of Picard iterations
  (\code{"picard"} only).
  \item[\code{picard_rtol} (float, $10^{-8}$)] Relative $L^2$ tolerance on
  the change in $\hne$ between Picard iterations (\code{"picard"} only).
  \item[{\code{picard_relax} (float in $(0,1]$, 1.0)}] Under-relaxation
  factor for the frozen KBM coefficient update:
  $c_{\mathrm{new}} = \text{relax}\cdot c_{\mathrm{computed}} + (1-\text{relax})\,c_{\mathrm{old}}$.
  Use $<1$ if the gate flip-flops between iterations; 1.0 disables it.
\end{description}
After a \code{"picard"} solve, \code{self.picard_info} records the
iteration history ($\bar\alpha$, gate state, profile change) and whether
the loop converged. \code{self.kbm_info} records the treatment settings.

\paragraph{Neutrals treatment.} \code{neutrals_treatment} controls how the
FC / CX neutral equations are solved.
\begin{itemize}
  \item \code{"fem"} (default): $n_{\mathrm{FC}}$ and $n_{\mathrm{CX}}$ are
  solved as part of the mixed three-field Firedrake system (legacy
  behaviour). Requires the FE mesh to resolve the neutral penetration depth
  $\lambda = |\VFC|/(n_e(S_i+\SCX))$; on neutral-opaque devices (e.g.\
  SPARC, $\lambda\sim0.2$\,mm) an unresolved layer produces oscillatory /
  negative neutral densities and Newton line-search failures.
  \item \code{"analytic"} (requires \code{kbm_treatment="picard"}): each
  Picard iteration, $n_{\mathrm{FC}}$ and $n_{\mathrm{CX}}$ are solved
  \emph{exactly} (integrating factor / per-step exponential integrator) on a
  dense \code{n_neutral_sub}-point sub-grid for the current $n_e$ (see
  \code{_solve_neutrals_analytic_nondim}), then frozen in the $n_e$
  equation (with the same under-relaxation as the KBM coefficients). The
  mixed-system residuals for $n_{\mathrm{FC}}$ / $n_{\mathrm{CX}}$ are
  replaced by trivial pinning forms so the rest of the machinery (mixed
  space, BCs, output extraction) is unchanged. Positivity of the neutrals
  is guaranteed and \code{x_res} only needs to resolve the \emph{electron}
  profile.
\end{itemize}
\code{n_neutral_sub} (int, 4001): number of sub-grid points for the
analytic neutral solve; choose so that $L/\code{n_neutral_sub}\ll\min\lambda$.

\paragraph{\code{"outer"} mode inside the solve.} One nonlinear solve is
the bare SNES solve in \code{"inner"} mode; in \code{"outer"} mode it is
the secant iteration that pins $\hne'(0)$ by treating the (unconstrained)
inner flux as the unknown. Inside the Picard loop the slope carries over
between iterations, so later shootings start already close. The selected
inner slope is reported as \code{self.dne_dx_inner_solved}.

\paragraph{Sets.} SI-unit attributes (parent-class compatibility):
\code{x_sol}, \code{ne_sol}, \code{nFC_sol}, \code{nCX_sol}, \code{u_fd}
(mixed Firedrake Function, storing hat values), \code{W_fd}, \code{V_fd}.
Non-dimensional attributes (diagnostics): \code{hat_x_sol},
\code{hat_ne_sol}, \code{hat_nFC_sol}, \code{hat_nCX_sol}, and the scales
\code{_L_nd}, \code{_n0_nd}, \code{_T0_nd}, \code{_S0_nd}, \code{_tau_nd},
\code{_V0_nd}, \code{_D0_nd}.

\paragraph{Returns.} The common solver result schema, in SI units.

\subsection{Module \code{__getattr__(name)}}
Resolves the pre-refactor class name \code{saarelma_connor_nondim} to the
assembled class \code{saarelma_connor}. Deferred rather than a top-level
import: \code{saarelma_connor_api} imports this module, so binding the name
at import time would be a cycle.

\end{document}
